\chapter{Discussion et limites}
\label{ch:disc}

\section{Limites de l'IA générative}
Malgré les garde-fous du RAG, les LLMs peuvent produire des
\emph{hallucinations} sur des questions très spécifiques absentes
du contexte. Le mode \og refus poli \fg{} hors sujet (cuisine,
météo, sport, etc.) a été ajouté en sprint~6 pour cantonner POSTIE
à son périmètre PO et adjacents (technique, communication d'équipe,
documentation). Cette restriction garantit un comportement
professionnel mais n'élimine pas le risque de réponse erronée à
l'intérieur du périmètre métier.

\section{Coût opérationnel et dépendance API}
La dépendance aux API OpenAI et Mistral entraîne un coût variable
proportionnel au volume d'utilisation. Au-delà du coût direct, cette
dépendance pose la question de la \emph{vendor lock-in} et de la
continuité de service en cas de panne ou de changement tarifaire.
Une étude d'alternatives auto-hébergées (Mistral 7B, Llama 3) est
prévue en perspective.

\section{Enjeux de confidentialité}
L'envoi de documents internes à des LLMs hébergés impose une
vigilance RGPD~: \textit{redaction} des PII avant ingestion,
signature de contrats DPA (\textit{Data Processing Agreement}) avec
les fournisseurs, ou bascule sur une alternative on-premise. POSTIE
documente explicitement ces enjeux pour l'utilisateur final dans
sa fiche produit.

\section{Limites méthodologiques}
L'évaluation a été menée sur un panel restreint de trois PO et sur
une période courte de deux semaines. Une étude longitudinale en
équipe réelle (sur un trimestre, avec mesure de la vélocité avant
et après adoption) serait nécessaire pour mesurer l'impact concret
sur la productivité et la qualité du backlog.

\section{Limites techniques résiduelles}
\begin{itemize}
  \item Le free tier Render impose un \emph{cold start} de
        $\sim$30~secondes après 15~minutes d'inactivité.
  \item L'absence de disque persistant en free tier impose une
        re-ingestion documentaire après chaque redeploy
        (résolu en sprint~6 par copie du backlog initial dans
        l'image Docker).
  \item Le routage multi-modèles est basé sur des heuristiques de
        mots-clés~; un classifieur entraîné apporterait
        probablement de meilleurs résultats.
\end{itemize}
